\documentclass[10pt,a4paper]{amsart}
\usepackage[margin=19mm]{geometry}
\usepackage{amsmath,amssymb,mathtools}
\usepackage[scr=rsfs,cal=euler]{mathalfa}
\usepackage{xfrac}
\usepackage[unicode,hidelinks,bookmarks=false]{hyperref}
\newcommand{\ie}{i.e.\ }
\newcommand{\cf}{cf.\ }
\newcommand{\resp}{respectively\ }
\newcommand{\Subsection}{\subsection}
\providecommand{\nobreakdash}{\nobreak\hbox{-}}
\setlength{\emergencystretch}{2em}
\multlinegap0pt
\usepackage{amssymb}
\usepackage{mathtools}
\usepackage{dsfont}
\usepackage{enumerate}
\usepackage{stackengine}
\newtheorem{proposition}{Proposition}
\newtheorem{theorem}{Theorem}
\newcommand{\Int}{\int_{\mathbb{R}^n}}
\newcommand{\N}{\mathbb{N}}
\newcommand{\R}{\mathbb{R}}
\newcommand{\bs}{\Bar{s}}
\renewcommand{\ge}{\geqslant}
\renewcommand{\le}{\leqslant}
\renewcommand{\leq}{\leqslant}
\title{Fredholm alternative for a general class\\ of nonlocal operators}
\author{Francesco De Pas, Serena Dipierro, Enrico Valdinoci}
\date{Source: arXiv:2604.07708v1 (CC BY 4.0)}
\begin{document}
\maketitle

\noindent\emph{Source and licence.} Fredholm alternative for a general class\\ of nonlocal operators. Version arXiv:2604.07708v1 (CC BY 4.0), distributed by arXiv under the Creative Commons Attribution 4.0 International licence, http://creativecommons.org/licenses/by/4.0/, as recorded in the arXiv licence metadata for this exact version and shown on the abstract page https://arxiv.org/abs/2604.07708v1. Full licence notice: \texttt{licenses/arxiv-2604.07708-CC-BY-4.0.txt}.\par\smallskip

\noindent\textit{Editorial scope.} Counted unit: the article's proposition crizzo (source0.tex lines 1487--1496). The article's complete original proof of it is delivered with it (source0.tex lines 1498--1649, one \texttt{proof} environment of 152 lines). No mathematics is written, repaired or re-derived: the statement and the proof are the article's own lines, in the article's own notation, and no retained line is substituted or re-flowed.  Retained uncounted as the source-local context the proof needs: lines 543--544; lines 707--711; lines 920--922; lines 929--931; lines 1132--1137; lines 1327--1329; lines 1331--1338.  Nothing was omitted from any retained span, so the package carries no omission record.  No retained line contains a citation, so the wrapper carries no bibliography and no bibliography entry.  No retained line contains a figure environment, an image-inclusion command or drawing code, so no image is delivered and the package declares no asset dependency.  Editorial formatting only, in addition: an AMS \texttt{amsart} preamble that restates the article's own declarations and macros used by the retained lines, namely the environments \texttt{proposition}, \texttt{theorem} on separate counters, together with the standard packages those lines need.  The retained environments print in this document as Proposition 1, Theorem 1 in order of appearance, whereas the article prints the same items with the numbering of its own full document.  The complete native source of the article as distributed in the arXiv e-print for this exact version is bundled unchanged as \texttt{sources/198110/source0.tex}.
			\begin{equation}\label{constdef}
\sup_{s \in [-1,1)} \frac{c_s}{(1-s)} < + \infty  \end{equation} and

			\begin{proposition}\label{usasempre}
				Let~$s\in (0,1)$,~$u \in \mathcal{D}(\R^n)$ and  let~$ R >0$ be such that~$B_{ R}$ contains the support of~$u$.			
				Then, for any~$x\in\R^n\setminus B_{ 2 R}$,
				\[  |D^s u (x)| \leq  \frac{  2^{n+s} c_s  \Vert u \Vert_{L^1\left(B_{ R}\right)}} {|x|^{n+s}} . \]
			\end{proposition}

			\begin{equation}\label{riesz}
				I_{\alpha}u (x) := K_{\alpha} \ast u (x) = \frac{1}{\gamma_{\alpha,n}}\Int \frac{u(y)}{|x-y|^{n-\alpha}} \, dy,   
			\end{equation}

			\begin{equation}\label{relation}
				c_s= \frac{n+s-1}{\gamma_{1-s}}.
			\end{equation}

			\begin{theorem}\label{usacchialo}
				Let~$\Bar{s} \in (0,1)$ and~$u \in \mathcal{D}(\R^n)$. Then, for any~$s \in (\bs,1]$,
				\begin{equation}\label{imp}
					D^{\bs} u= I_{s-\bs} D^s u.
				\end{equation}
			\end{theorem}

		\begin{equation}\label{simn}
			\Vert u \Vert_{L^p(\Omega)} \leq \frac{C}{s} \Vert D^s u \Vert_{L^p(\R^n)}.
		\end{equation}

		\begin{proposition}\label{costa}
			Let~$s \in (0,1)$ and~$\Omega$ be a bounded domain of~$\R^n$. Let~$\rho>0$ be such that~${\Omega \subset [-\rho,\rho]^n}$.
			
			Then, there exists~$C>0$, depending only on~$n$, such that, for any~$u \in H_0^{s,1}(\Omega)$,
			\begin{equation*}
				\Vert u \Vert_{L^1(\Omega)} \leq \frac{C\,\rho^s }{s} \Vert D^s u \Vert_{L^1(\R^n)}.
			\end{equation*}
		\end{proposition}

		\begin{proposition}\label{crizzo}
		Let~$\bs \in (0,1)$ and~$\Omega$ be a bounded domain of~$\R^n$.
			
			 Then, there exists~$C>0$,
			 depending only on~$n$, $\bs$ and~$\Omega$,
			 such that, for any~$s \in [\bs,1)$, $p \in [1,+\infty)$ and~$u \in H_0^{s,p}(\Omega)$, 
			\begin{equation}\label{re}
				\Vert D^{\bs}u \Vert_{L^p(\R^n)} \leq C\Vert D^s u \Vert_{L^p(\R^n)}. 
			\end{equation}
		\end{proposition}

		\begin{proof}
If~$s = \bar{s}$, the desired result holds true by taking~$C=1$, thus from now on
we suppose that~$s\in (\bar{s},1)$. 
In this case, we first establish Proposition~\ref{crizzo} for~$u \in \mathcal{D}(\Omega)$ and then we apply a density argument to complete the proof. 
						
			We take~$R>0$ such that~$\Omega \subset B_{R}$. Also, recalling~\eqref{relation}, we see that
			\[ \gamma_{s-\bs} = \frac{n-(s- \bs)}{c_{1-(s-\bs)}},\]
and therefore, in light of~\eqref{constdef},			
			\begin{equation}\label{fgh} \widetilde\gamma:= \sup_{ s \in (\bs,1)} \frac{n-(s- \bs)}{\gamma_{s-\bs}(s-\bs)} =
				\sup_{ s \in (\bs,1)} \frac{c_{1-(s-\bs)}}{s-\bs} < +\infty. 
				\end{equation}
					
	Now we prove the desired result for~$p=1$.
To this end, we define
	\begin{align*}
	 &A_I := \frac{1}{\gamma_{s-\bs}} \int_{B_{2R}} \left( \int_{B_{4R}} \frac{|D^su(y)|}{|x-y|^{n-(s-\bs)}} \, dy \right) \, dx, \\
	  &A_{II} := \frac{1}{\gamma_{s-\bs}} \int_{B_{2R}} \left( \int_{\R^n \setminus B_{4R}} \frac{|D^su(y)|}{|x-y|^{n-(s-\bs)}} \, dy \right) \, dx \\ {\mbox{and }}\qquad
	 &B := \int_{\R^n \setminus B_{2R}} |D^{\bar{s}} u (x)| \, dx.
	\end{align*}
	We first estimate~$A_I$: making use of~\eqref{fgh}, we find that 			
\begin{equation}\label{jeiwoytgrej768905}\begin{split}
A_I &\leq   \frac{1}{\gamma_{s-\bs}}  \left( \int_{B_{6R}} \frac{dz}{|z|^{n-(s-\bs)}} \right)\Vert D^s u \Vert_{L^1(\R^n)}\\&
= \frac{S_{n-1}(6R)^{s- \bar{s}}}{\gamma_{s-\bs}(s-\bar{s})} \Vert D^s u \Vert_{L^1(\R^n)}\\
&\le  \frac{6S_{n-1}\widetilde\gamma\max\{1,R\} }{n-(s-\bs)} \Vert D^s u \Vert_{L^1(\R^n)}\\&
\le  \frac{C_1 }{\bs} \Vert D^s u \Vert_{L^1(\R^n)},
\end{split}\end{equation} for some~$C_1>0$, depending only on~$n$ and~$\Omega$.

Now we estimate~$A_{II}$. 
To this end, we set~$\Tilde{c}:= \sup_{s \in(-1,1)} c_{s}$, which is finite, thanks to~\eqref{constdef}.
As a result, 
exploiting Propositions~\ref{usasempre} and~\ref{costa},
\begin{equation}\label{jeiwoytgrej7689052}\begin{split}
			A_{II}& \leq \frac{ 2 ^{n+s}\,\Tilde{c}}{\gamma_{s-\bar{s}}} \left[ \int_{\R^n \setminus B_{4R}} \frac{1}{|y|^{n+s}} \left( \int_{B_{2R}} \frac{dx}{|x-y|^{n-(s-\bar{s})}} \right) \, dy \right] \Vert u \Vert_{L^1(\Omega)}	 \\ 
& \leq \frac{ 2^{3n+\bs}\,\Tilde{c} \,S_{n-1} R^n}{\gamma_{s-\bar{s}}} \int_{\R^n \setminus B_{4R}} \frac{dy}{|y|^{2n + \bar{s}}}\,\Vert u \Vert_{L^1(\Omega)}  \\&
\le\frac{ 2^{n}\,\Tilde{c} \,S^2_{n-1}}{\gamma_{s-\bar{s}} \, R^{\bs} } \,\Vert u \Vert_{L^1(\Omega)} \\
& \le 
\frac{ 2^{n}\,\Tilde{c} \,\widetilde\gamma\,S^2_{n-1}(s-\bs)}{(n-(s-\bs)) R^{\bs} } \,\Vert u \Vert_{L^1(\Omega)}\\& \le
\frac{ 2^{n}\,\Tilde{c} \,\widetilde\gamma\,S^2_{n-1}}{\bs\,\min\{1, R\} } \,\Vert u \Vert_{L^1(\Omega)}\\&
\le\frac{ C_2}{\bar{s}^2} \Vert D^s u \Vert_{L^1(\R^n)},
			\end{split}\end{equation} for some~$ C_2>0$, depending
			only on~$n$ and~$\Omega$.

In order to estimate~$B$, we use Proposition~\ref{usasempre} and~\ref{costa}
to obtain that
\begin{equation}\label{jeiwoytgrej7689053}\begin{split}&
B \leq 2^{n+\bar{s}} \Tilde{c} \left( \int_{\R^n \setminus B_{2R}}\frac{dx}{|x|^{n+\bs}}\right) \Vert u \Vert_{L^1(\Omega)} 
\\&\qquad= \frac{ S_{n-1} 2^{n} \Tilde{c}}{\bar{s}\,R^{\bar{s}}}  \Vert u \Vert_{L^1(\Omega)}
\leq \frac{C_3}{\bar{s}^2} \Vert  D^s u \Vert_{L^1(\R^n)},\end{split}
\end{equation} for some~$C_3>0$, depending only on~$n$ and~$\Omega$.
			
Now we recall that, in light of Theorem~\ref{usacchialo}, $D^{\bar{s}}u =I_{s-\bar{s}} D^s u$, and thus, recalling the definition of~$I_{s-\bar{s}}$ 
in~\eqref{riesz}
\begin{eqnarray*}&& \Vert D^{\bar{s}}u \Vert_{L^1(\R^n)}=\int_{\R^n}| D^{\bar{s}}u(x)|\,dx
=\int_{\R^n}| K_{s-\bar s} \ast D^su(x)|\,dx\\&&\qquad\qquad
=\frac{1}{\gamma_{s-\bar s}}\int_{\R^n}\left|
\Int \frac{D^su(y)}{|x-y|^{n-(s-\bar s)}} \, dy\right|\,dx.\end{eqnarray*}
As a consequence of this and the estimates provided in~\eqref{jeiwoytgrej768905},
\eqref{jeiwoytgrej7689052} and~\eqref{jeiwoytgrej7689053}, we conclude that
$$
\Vert D^{\bar{s}}u \Vert_{L^1(\R^n)} \leq A_{I} + A_{II} + B\le\frac{C_4}{\bs}\left(1+\frac{1}{\bs}\right) \Vert  D^s u \Vert_{L^1(\R^n)},
$$ for some~$C_4>0$, depending only on~$n$ and~$\Omega$.
Hence, the proof of the desired claim
is complete for~$p=1$ and~$u \in \mathcal{D}(\Omega)$.
			
			We now establish the desired result for~$p \in (1, +\infty)$.
			For this, we define
			\begin{align*}
			&A_{I}(p):= \frac{2^p}{\gamma^p_{s-\bs}} \int_{B_{2R}} \left( \int_{B_{4R}} \frac{|D^su(y)|}{|x-y|^{n-(s-\bs)}} \, dy\right)^p \, dx,  \\
			&A_{II}(p):= \frac{2^p}{\gamma^p_{s-\bs}} \int_{B_{2R}} \left (\int_{\R^n \setminus B_{4R}} \frac{|D^su(y)|}{|x-y|^{n-(s-\bs)}} \, dy\right)^p \, dx \\ {\mbox{and }} \qquad
			&B(p):= \int_{\R^n \setminus B_{2R}} |D^{\bs}u(x)|^p \, dx.
			\end{align*}
			We first estimate~$A_I(p)$. For this purpose, we observe that, if~$x \in B_{2R}$, setting~$p':=\frac{p}{p-1}$ and employing the H\"older inequality,
			\begin{eqnarray*}&&
				\int_{B_{4R}} \frac{|D^su(y)|}{|x-y|^{n-(s-\bs)}} \, dy\\& =&
				 \int_{B_{4R}} \frac{|D^su(y)|}{|x-y|^{\frac{n-(s-\bs)}{p}}} \frac{1}{|x-y|^{\frac{n-(s-\bs)}{p'}}} \, dy  \\ 
				 &\leq&  \left(\int_{B_{6R}} \frac{dz}{|z|^{n-(s-\bs)}} \right)^{\frac{1}{p'}} \left(\int_{B_{4R}} \frac{|D^s u(y)|^p}{|x-y|^{n-(s-\bs)}} \, dy \right)^{\frac{1}{p}}
				  \\ &= &\left( \frac{ S_{n-1} (6R)^{s-\bar{s}}}{s-\bs}\right)^{\frac{1}{p'}} \left(\int_{B_{4R}} \frac{|D^s u(y)|^p}{|x-y|^{n-(s-\bs)}} \, dy \right)^{\frac{1}{p}}  
\\ &\le &
\frac{ 6S_{n-1} \max\{1,R\} }{(s - \bar{s})^{\frac{1}{p'}}} \left(\int_{B_{4R}} \frac{|D^s u(y)|^p}{|x-y|^{n-(s-\bs)}} \, dy \right)^{\frac{1}{p}}. 
\end{eqnarray*}
Consequently, 
\begin{equation*}\begin{split}
A_I(p)& \leq \frac{(12 S_{n-1} \max\{1,R\})^p}{\gamma^p_{s-\bs}(s-\bs)^{\frac{p}{p'}}} \int_{B_{4R}} \left(\int_{B_{R_2}} \frac{dx}{|x-y|^{n-(s-\bs)}}  \right) |D^s u(y)|^p\, dy \\ &
=\frac{12^{p+s-\bar{s}} S_{n-1}^{p+2} \max\{1,R\}^p\,R^{s-\bar s}}{\gamma^p_{s-\bs}(s-\bs)^{\frac{p}{p'}+1}}  \Vert D^s u \Vert^p_{L^p(\R^n)}
\\&=\frac{12^{p+s-\bar{s}} S_{n-1}^{p+2} \max\{1,R\}^p\,R^{s-\bar s}}{\big(\gamma_{s-\bs}(s-\bs)\big)^p} \Vert D^s u \Vert^p_{L^p(\R^n)}\\&
\le\frac{12^{p+1} S_{n-1}^{p+2} \max\{1,R\}^{p+1}\,\widetilde\gamma^p}{\big(n-(s-\bar s)\big)^p} \Vert D^s u \Vert^p_{L^p(\R^n)}\\&\le
\frac{12^{p+1} S_{n-1}^{p+2} \max\{1,R\}^{p+1}\,\widetilde\gamma^p}{\bar{s}^p} \Vert D^s u \Vert^p_{L^p(\R^n)}
.
%%\le\frac{C_5}{\bar{s}^p} \Vert D^s u \Vert^p_{L^p(\R^n)},
			\end{split}\end{equation*}
		%%	for some~$C_5>0$, depending only on~$n$, $p$ and~$\Omega$.
			
			Now we estimate~$A_{II}(p)$. By Proposition~\ref{usasempre} and equation~\eqref{simn}, we see that
\begin{equation*}\begin{split}
A_{II}(p) &\leq \frac{2^{p+n+s} \Tilde{c}}{\gamma^p_{s-\bar{s}}}  \left[ \int_{B_{2R}} \left( \int_{\R^n \setminus B_{4R}} \frac{dy}{|y|^{n+s} |x-y|^{n-(s-\bs)}}\right)^p \, dx \right] \Vert u \Vert_{L^1(\Omega)}^p \\ & 
\leq  \frac{2^{p+n+s+p(n-s+\bar s)} \Tilde{c}}{\gamma^p_{s-\bar{s}}}
\left(\int_{\R^n \setminus B_{4R}} \frac{dy}{|y|^{2n+ \bs}}\right)^p \Vert u \Vert_{L^1(\Omega)}^p \\&\le
\frac{2^{p+n+s+p(n-s+\bar s)} \Tilde{c}\,S_{n-1}^p|\Omega|^{p-1}}{(4R)^{p(n+\bar s)}\gamma^p_{s-\bar{s}}} \Vert u \Vert_{L^p(\Omega)}^p\\&  \leq
\frac{2^{p+n+s+p(n-s+\bar s)} \Tilde{c}\,S_{n-1}^p|\Omega|^{p-1}\widetilde\gamma^p}{(4R)^{p(n+\bar s)}\bar{s}^p} \Vert u \Vert_{L^p(\Omega)}^p\\&\le 
\frac{C\,2^{p+n+s+p(n-s+\bar s)} \Tilde{c}\,S_{n-1}^p|\Omega|^{p-1}\widetilde\gamma^p}{(4R)^{p(n+\bar s)}\bar{s}^{2p}} \Vert D^su \Vert_{L^p(\Omega)}^p
,\end{split}\end{equation*}
where~$C$ depends only on~$n$ and~$\Omega$.
			
Now we estimate~$B(p)$. We will proceed in a similar way as for~$p=1$,
using Proposition~\ref{usasempre} and equation~\eqref{simn}. The details go as follows:
\begin{equation*}\begin{split}
B(p)& \leq (2^{n+ \bar{s}} \Tilde{c})^p  \left( \int_{\R^n \setminus B_{2R}} \frac{dx}{|x|^{p(n+ \bs)}}\right) \Vert u \Vert^p_{L^1(\Omega)} \\
& =  \frac{(2^{n+ \bar{s}} \Tilde{c})^p S_{n-1} }{(n(p-1)+ \bar{s}p) R^{n(p-1)+ \bar{s}p}} \Vert u \Vert^p_{L^1(\Omega)}\\&\leq 
\frac{(2^{n+ \bar{s}} \Tilde{c})^p S_{n-1}\,|\Omega|^{p-1} }{(n(p-1)+ \bar{s}p) R^{n(p-1)+ \bar{s}p}} \Vert u \Vert^p_{L^p(\Omega)}
\\& \le \frac{C(2^{n+ \bar{s}} \Tilde{c})^p S_{n-1}\,|\Omega|^{p-1} }{(n(p-1)+ \bar{s}p) R^{n(p-1)+ \bar{s}p}\bar{s}^p}  \Vert D^s u \Vert^p_{L^p(\R)}.
\end{split}\end{equation*}

Gathering these estimates and recalling
Theorem~\ref{usacchialo}, we conclude that
\begin{equation}\label{eiwoty437683869wejfo}\begin{split}&
\Vert D^{\bs} u \Vert^p_{L^p(\R^n)} \le A_{I}(p) + A_{II}(p) + B(p)\\
&\qquad \le\left(
\frac{12^{p+1} S_{n-1}^{p+2} \max\{1,R\}^{p+1}\,\widetilde\gamma^p}{\bar{s}^p}\right.\\&\qquad\qquad\qquad
+\frac{C\,2^{p+n+s+p(n-s+\bar s)} \Tilde{c}\,S_{n-1}^p|\Omega|^{p-1}\widetilde\gamma^p}{(4R)^{p(n+\bar s)}\bar{s}^{2p}}\\ &\qquad\qquad\qquad \left. 
+\frac{C(2^{n+ \bar{s}} \Tilde{c})^p S_{n-1}\,|\Omega|^{p-1} }{(n(p-1)+ \bar{s}p) R^{n(p-1)+ \bar{s}p}\bar{s}^p} \right) \Vert D^s u \Vert^p_{L^p(\R)}.
\end{split}\end{equation}

Now we observe that, for all~$a$, $b\ge0$, we have that~$a^p+b^p\le(a+b)^p$, therefore,
for every~$j\in\N$ and~$a_1,\dots, a_j\ge0$, we have that
$$ \sum_{i=1}^j a_i^p\le \left( \sum_{i=1}^j a_i\right)^p.$$
Using this into~\eqref{eiwoty437683869wejfo}, we thereby obtain that
\begin{equation}\label{dewiosdvjt473kipipoiupuoy}\begin{split}
\Vert D^{\bs} u \Vert_{L^p(\R^n)}& \le\left(
\frac{12^{p+1} S_{n-1}^{p+2} \max\{1,R\}^{p+1}\,\widetilde\gamma^p}{\bar{s}^p}\right.\\ &\qquad \qquad
+\frac{C\,2^{p+n+s+p(n-s+\bar s)} \Tilde{c}\,S_{n-1}^p|\Omega|^{p-1}\widetilde\gamma^p}{(4R)^{p(n+\bar s)}\bar{s}^{2p}}\\ &\qquad\qquad \left. 
+\frac{C(2^{n+ \bar{s}} \Tilde{c})^p S_{n-1}\,|\Omega|^{p-1} }{(n(p-1)+ \bar{s}p) R^{n(p-1)+ \bar{s}p}\bar{s}^p} \right)^{\frac1p} \Vert D^s u \Vert_{L^p(\R)}\\
&\le\left(
\frac{12^{1+\frac1p} S_{n-1}^{1+\frac2p} \max\{1,R\}^{1+\frac1p}\,\widetilde\gamma}{\bar{s}}\right.\\ &\qquad \qquad
+\frac{C^{\frac1p}\,2^{1+n-s+\bar s+\frac{n+s}p} \Tilde{c}^{\frac1p}\,S_{n-1}|\Omega|^{\frac{p-1}p}\widetilde\gamma}{(4R)^{n+\bar s}\bar{s}^{2}}\\ &\qquad \qquad \left. 
+\frac{C^{\frac1p} 2^{n+ \bar{s}} \Tilde{c}\, S_{n-1}^{\frac1p}\,|\Omega|^{\frac{p-1}p} }{(n(p-1)+ \bar{s}p)^{\frac1p} R^{\frac{n(p-1)+ \bar{s}p}p}\bar{s}} \right) \Vert D^s u \Vert_{L^p(\R)}
.\end{split}\end{equation}

Moreover, we point out that, for all~$M\ge0$, we have that
$$\min\{1,M\}\le M^{\frac1p}\le \max\{1,M\}.$$
{F}rom this and~\eqref{dewiosdvjt473kipipoiupuoy}, we infer that
$$ \Vert D^{\bs} u \Vert_{L^p(\R^n)}\le \widetilde C\Vert D^s u \Vert_{L^p(\R)},$$
for some~$\widetilde C>0$, depending only on~$n$, $\bs$ and~$\Omega$, as desired.
					
{N}ow we complete the proof Proposition~\ref{crizzo} by
employing a density argument. More precisely, let~$u \in H^{s,p}_0(\Omega)$ and~$(\phi_k)\in\mathcal{D}(\Omega)$ converging to~$u$ in~$u \in H^{s,p}_0(\Omega)$ as~$k\to+\infty$. Then, for all~$k\in\N$,
\begin{equation}\label{nevren}
				\Vert D^{\bs}\phi_k \Vert_{L^p(\R^n)} \leq C \Vert D^s \phi_k \Vert_{L^p(\R^n)}. 
\end{equation}
Furthermore, since~$(D^s \phi_k)$ is a Cauchy sequence in $L^p(\R^n, \R^n)$, we see that $(D^{\bar{s}}\phi_k)$ is a Cauchy sequence in $L^p(\R^n, \R^n)$ too,
thanks to~\eqref{nevren}.
As a result, the thesis plainly follows passing to the limit as $k \to +\infty$ in~\eqref{nevren}.
\end{proof}

\end{document}
